\section{Continuum Limit Construction: The Program}
\label{sec:continuum-program}

\textbf{Status: [PROGRAM] / [RESEARCH ROADMAP]}

The construction of the continuum limit requires controlling the theory as the lattice spacing $a \to 0$ and the coupling $\beta \to \infty$. This is the domain of Constructive Quantum Field Theory.

\subsection{The Target Theorem: Existence of the Continuum Limit}

\begin{conjecture}[Existence of Continuum Limit]
\label{conj:continuum-limit}
There exists a renormalization group trajectory $\beta(a)$ such that the limit of the Schwinger functions exists:
\[
S_n(x_1, \dots, x_n) = \lim_{a \to 0} Z_a^{-n/2} \langle \phi_a(x_1) \dots \phi_a(x_n) \rangle_{\beta(a)}
\]
where $\phi_a$ are suitable lattice fields and $Z_a$ is a wavefunction renormalization constant.
Moreover, the limiting functions $S_n$ satisfy the Osterwalder-Schrader axioms, allowing for the reconstruction of a relativistic QFT on Minkowski space.
\end{conjecture}

\subsection{Proposed Method: Balaban's Renormalization Group}

The most promising route to proving Conjecture \ref{conj:continuum-limit} is the rigorous Renormalization Group method developed by Tadeusz Balaban.

\begin{itemize}
    \item \textbf{Multiscale Analysis}: The field is decomposed into fluctuation fields at different scales.
    \item \textbf{Effective Actions}: One proves that the effective action at each scale remains close to a local action (the "relevant" part) plus small "irrelevant" corrections.
    \item \textbf{Stability}: The key difficulty is proving stability of the flow, i.e., that the effective actions do not diverge.
\end{itemize}

\subsection{Proof Obligations}

To complete this step, one must:
\begin{enumerate}
    \item \textbf{Verify Hypotheses}: Explicitly check that the $SU(N)$ Wilson action satisfies the initial conditions required by Balaban's theorems.
    \item \textbf{Adapt to Finite Volume}: Ensure the bounds hold uniformly in the volume $L$ to allow for the thermodynamic limit.
    \item \textbf{OS Axioms}: Verify that the limiting functions satisfy Reflection Positivity and Cluster Decomposition.
\end{enumerate}

This program, while advanced, provides a concrete path to the continuum limit.
